Esta granularidad no es habitual en soluciones generalistas,
que suelen tratar todo como subida de fichero sin reglas contextuales.

\subsubsection*{Versionado y reenv\'io con trazabilidad}

Cada entrega queda registrada con versi\'on y fecha, manteniendo un
historial completo y una versi\'on activa por entregable.
Cuando el reenv\'io est\'a permitido, el sistema conserva la trazabilidad
sin sobrescribir el historial previo, facilitando auditor\'ia docente.

\subsubsection*{Operativa docente optimizada}

Se incorporan descargas masivas por entregable o actividad,
previsualizaci\'on de archivos en l\'inea y listado de contenido
de ZIP, reduciendo tareas manuales de comprobaci\'on y clasificaci\'on.

\subsubsection*{Robustez en el flujo de estudiante}

El flujo de entrega incorpora borradores autom\'aticos en cliente,
revalidaci\'on antes de enviar y mensajes de error contextualizados.
Esto minimiza p\'erdida de trabajo en conexiones inestables
y mejora la experiencia del alumnado.

\section{Alternativas de arquitectura y stack}

\subsection{Arquitectura de aplicaci\'on}

Se valoraron dos enfoques principales:

\begin{itemize}
	\item \textbf{A1 - SPA + API desacoplada} (opci\'on elegida).
	\item \textbf{A2 - Aplicaci\'on monol\'itica server-side renderizada}.
\end{itemize}

\textbf{A1} permite separar claramente la capa de experiencia de usuario
de la l\'ogica de negocio y persistencia, facilitando pruebas independientes,
despliegues diferenciados y evoluci\'on incremental.

\textbf{A2} reduce complejidad inicial de infraestructura,
pero limita flexibilidad de frontend y complica evoluci\'on de UX avanzada
en escenarios con roles y paneles diferenciados.

\subsection{Backend framework}

Se compararon principalmente:

\begin{itemize}
	\item \textbf{B1 - Spring Boot (Java)} (elegida).
	\item \textbf{B2 - Node.js (Express/Nest)}.
	\item \textbf{B3 - Django (Python)}.
\end{itemize}

La elecci\'on de \textbf{Spring Boot} se justifica por:

\begin{itemize}
	\item ecosistema maduro para seguridad, JPA, validaci\'on y testing,
	\item integraci\'on natural con migraciones Flyway,
	\item buen soporte para arquitectura por capas y control transaccional.
\end{itemize}

Node.js y Django son opciones v\'alidas,
pero para este proyecto requer\'ian m\'as decisiones de ensamblaje
en piezas de seguridad/persistencia para alcanzar un nivel equivalente
en el tiempo disponible.

\subsection{Frontend framework}

Alternativas consideradas:

\begin{itemize}
	\item \textbf{F1 - React + TypeScript + Vite} (elegida).
	\item \textbf{F2 - Angular}.
	\item \textbf{F3 - Vue}.
\end{itemize}

React con TypeScript y Vite ofrece buen equilibrio entre velocidad de desarrollo,
ecosistema y control de complejidad para una SPA de tama\~no medio.
Angular aporta una estructura m\'as opinionada,
pero con curva de entrada superior para el alcance temporal del TFG.
Vue mantiene una DX excelente, aunque en este caso la experiencia previa del equipo
y la alineaci\'on con librer\'ias de testing/CI favorecieron React.

\section{Alternativas de persistencia y almacenamiento}

\subsection{Base de datos relacional}

Opciones comparadas:

\begin{itemize}
	\item \textbf{D1 - PostgreSQL gestionado} (elegida en producci\'on).
	\item \textbf{D2 - MySQL gestionado}.
	\item \textbf{D3 - H2 embebida} (solo desarrollo/pruebas locales).
\end{itemize}

PostgreSQL se mantuvo como opci\'on principal por estabilidad, compatibilidad
con el stack y buen comportamiento en entornos PaaS.
H2 se reserva para contextos de desarrollo y pruebas, no para producci\'on.

Adem\'as, se produjo una migraci\'on de proveedor de Postgres (Neon a Render Postgres)
sin cambios de c\'odigo de dominio,
apoy\'andose en variables de entorno y configuraci\'on portable del datasource,
lo que confirma bajo acoplamiento a proveedor concreto.

\subsection{Estrategia de esquema}

Se compararon:

\begin{itemize}
	\item \textbf{E1 - Flyway como autoridad de esquema} (elegida).
	\item \textbf{E2 - Hibernate \texttt{ddl-auto=update} en producci\'on}.
\end{itemize}

Se adopta \textbf{E1} por trazabilidad y control expl\'icito de cambios,
reduciendo riesgo de divergencias de esquema entre entornos.

\subsection{Almacenamiento de archivos}

Alternativas evaluadas:

\begin{itemize}
	\item \textbf{S1 - Almacenamiento local en servidor}.
	\item \textbf{S2 - Cloudinary} (elegida para producci\'on persistente).
	\item \textbf{S3 - OneDrive/Microsoft Graph} (integraci\'on opcional por contexto acad\'emico).
\end{itemize}

La opci\'on local es simple y adecuada en desarrollo,
pero en plataformas con sistema de archivos ef\'imero (como despliegues free)
puede provocar p\'erdida de datos tras reinicio.
Por ello se adopta \textbf{Cloudinary} en producci\'on.

La integraci\'on \textbf{OneDrive} aporta valor en casos docentes
que requieren sincronizaci\'on con ecosistemas Microsoft,
pero se mantiene como capacidad configurable y no obligatoria del n\'ucleo.

\section{Alternativas de seguridad y sesi\'on}

\subsection{Modelo de autenticaci\'on}

Se compararon dos enfoques principales:

\begin{itemize}
	\item \textbf{A1 - JWT (access + refresh)} (elegido).
	\item \textbf{A2 - Sesi\'on server-side tradicional}.
\end{itemize}

En una SPA desacoplada, JWT con refresh ofrece una experiencia m\'as natural
para cliente API-first y facilita escalado stateless del backend.
La sesi\'on server-side reduce cierta complejidad de token,
pero exige estrategia adicional de estado distribuido en despliegues horizontales.

\subsection{Extensiones de seguridad}

Se han considerado funcionalidades avanzadas (inicio de sesi\'on federado con OAuth2/OIDC y 2FA)
como l\'inea de evoluci\'on.
La decisi\'on de no imponerlas en el MVP responde al criterio
de maximizar cobertura funcional core sin comprometer plazos.

\section{Alternativas investigadas y descartadas}

\subsection{SharePoint para OneDrive corporativo}

Se valor\'o el uso de SharePoint como intermediario para acceder a cuentas
corporativas de OneDrive. La opci\'on facilitaba el encaje institucional,
pero implicaba perder capacidades clave que s\'i ofrece la integraci\'on directa,
como la selecci\'on de carpeta de destino o la gesti\'on flexible de rutas
durante la subida de entregas. Dado que estas funciones son parte del valor
diferencial de la plataforma, se descart\'o esta v\'ia.

\subsection{Registro de usuarios con OAuth}

Se estudi\'o habilitar registro de usuarios mediante OAuth/OIDC
para que el alumnado pudiera crear cuentas de forma aut\'onoma.
Finalmente se prioriz\'o el flujo habitual en entornos acad\'emicos,
donde el administrador crea usuarios y los asigna a cursos y grupos.
Esto evita altas desalineadas con la matr\'icula real y simplifica el control
de permisos y pertenencia a grupos.

\section{Alternativas de despliegue y operaci\'on}

\subsection{Estrategia de despliegue}

Opciones principales:

\begin{itemize}
	\item \textbf{P1 - PaaS con automatizaci\'on GitHub Actions + Render} (elegida).
	\item \textbf{P2 - VPS autogestionado con pipeline manual o Jenkins}.
	\item \textbf{P3 - Orquestaci\'on completa en Kubernetes}.
\end{itemize}

\textbf{P1} reduce carga operativa,
permite despliegue r\'apido y mantiene un coste ajustado para contexto acad\'emico.
\textbf{P2} y \textbf{P3} ofrecen m\'as control,
pero su complejidad no aporta retorno proporcional dentro del alcance del TFG.

\subsection{CI/CD y calidad}

Se compararon:

\begin{itemize}
	\item \textbf{C1 - CI/CD integrada en GitHub Actions} (elegida).
	\item \textbf{C2 - CI externa (GitLab CI/Jenkins) con mayor administraci\'on}.
\end{itemize}

La alternativa elegida integra lint, tests, build, an\'alisis de seguridad,
mutaci\'on y publicaci\'on de artefactos en un flujo unificado y reproducible.

\section{Comparativa multicriterio}

La Tabla~\ref{tab:comparativa-global} sintetiza la comparaci\'on global
entre enfoques de soluci\'on.

\begin{table}[ht]
\centering
\small
\begin{tabular}{|P{3.9cm}|c|c|c|}
\hline
\textbf{Criterio} & \textbf{Soluci\'on elegida} & \textbf{Alternativa intermedia} & \textbf{Alternativa avanzada} \\
\hline
Tiempo de entrega MVP & 5 & 3 & 2 \\
\hline
Mantenibilidad & 4 & 3 & 4 \\
\hline
Coste operativo & 5 & 3 & 2 \\
\hline
Escalabilidad potencial & 4 & 3 & 5 \\
\hline
Complejidad de operaci\'on & 4 & 3 & 1 \\
\hline
Adecuaci\'on al TFG & 5 & 3 & 2 \\
\hline
\end{tabular}
\caption{Comparativa global (escala 1-5, mayor es mejor)}
\label{tab:comparativa-global}
\end{table}

Interpretaci\'on:
la alternativa seleccionada no maximiza todos los ejes te\'oricos,
pero s\'i ofrece el mejor equilibrio entre valor entregado,
riesgo t\'ecnico y viabilidad temporal del proyecto.

\section{Decisiones adoptadas y justificaci\'on final}

Las decisiones finales que mejor explican el resultado del proyecto son:

\begin{itemize}
	\item arquitectura SPA + API desacoplada para separar responsabilidades;
	\item backend Spring Boot con persistencia relacional y migraciones Flyway;
	\item frontend React + TypeScript para productividad con control tipado;
	\item almacenamiento cloud persistente en producci\'on y local en desarrollo;
	\item despliegue en PaaS con CI/CD automatizada y verificaciones de calidad.
\end{itemize}

Esta combinaci\'on ha permitido llegar a un sistema funcional,
con cobertura de pruebas alta y capacidad de evoluci\'on,
sin sobredimensionar infraestructura ni desviarse del objetivo acad\'emico.

\section{Conclusiones del cap\'itulo}

La comparaci\'on confirma que la soluci\'on implementada no responde a una selecci\'on
arbitraria de herramientas,
sino a una evaluaci\'on razonada frente a alternativas reales.

El enfoque adoptado prioriza equilibrio entre robustez, coste y mantenibilidad,
lo que lo hace especialmente adecuado para un TFG de perfil profesional
con restricciones de tiempo y recursos.
	
	 


























